\ifx\litemode\undefined\documentclass{article}\fi
\usepackage[margin=3cm]{geometry}

\usepackage{graphicx}

\usepackage[table]{xcolor}
\usepackage{tikz}
\usetikzlibrary{shapes,arrows}
\definecolor{xlightgray}{HTML}{D1DBE4}
\definecolor{xgreen}{HTML}{B9C89F}
\definecolor{xred}{HTML}{F19F9E}
\definecolor{xyellow}{HTML}{F6DA71}
\definecolor{xblue}{HTML}{ACE7EF}
\definecolor{xpurple}{HTML}{C3A4E7}
\usepackage{todonotes}
\usepackage{booktabs,tabularx,makecell,hyperref}
\definecolor{rowgray}{RGB}{246,246,246}

\usepackage{pdflscape}
\usepackage{environ}

\usetikzlibrary{arrows.meta,positioning,fit,calc,decorations.pathmorphing}

% Math
\usepackage{amsmath,amssymb}
\usepackage{mathtools}
\usepackage{xspace}

% Algorithmic
\usepackage[noend]{algpseudocode}
\usepackage{amsthm}
\usepackage{cleveref}

\usepackage{float}
\floatstyle{ruled}
\newfloat{algo}{htbp}{algo}
\floatname{algo}{Algorithm}

% --- Custom commands ---
\newcommand{\genesis}{\ensuremath{B_{\mathrm{gen}}}}
\newcommand{\T}[0]{\ensuremath{\mathcal{T}}}
\newcommand{\slot}[0]{\ensuremath{\operatorname{\mathrm{slot}}}}

\newcommand{\target}{\mathsf{target}}
\newcommand{\kind}{\mathsf{kind}}
\newcommand{\height}{\mathsf{height}}
\newcommand{\validator}{\mathsf{validator}}
\newcommand{\just}{\ensuremath{\mathsf{justify}}\xspace}
\newcommand{\notar}{\ensuremath{\mathsf{notarize}}\xspace}
\newcommand{\parent}{\mathsf{parent}}

% Vote fields and state variables
\newcommand{\finalizeHeight}{\mathsf{finalize\_height}}
\newcommand{\finalizeTarget}{\mathsf{finalize\_target}}
\newcommand{\jtargets}{\mathsf{jtargets}}
\newcommand{\ntargets}{\mathsf{ntargets}}
\newcommand{\voteSlots}{\mathsf{voteSlots}}
\newcommand{\penalties}{\mathsf{penalties}}
\newcommand{\fresh}{\mathsf{fresh}}
\newcommand{\cnt}{\mathsf{count}}
\newcommand{\hash}{\mathsf{hash}}
\newcommand{\key}{\mathsf{key}}
\newcommand{\bit}{\mathsf{bit}}
\newcommand{\newKey}{\mathsf{newKey}}
\newcommand{\currentSlot}{\mathsf{currentSlot}}
\newcommand{\syncer}{\mathsf{syncer}}
\newcommand{\start}{\mathsf{start}}
\newcommand{\vote}{\mathsf{vote}}

% Common (across honest validators) marker
\newcommand{\comm}[1]{\bar{#1}}

% Horizontal separator rule
\newcommand{\sep}{\par\noindent\rule{\linewidth}{0.4pt}\par\medskip}

% Procedure names
\newcommand{\processSlot}{\textsf{processSlot}\xspace}
\newcommand{\processBlock}{\textsf{processBlock}\xspace}
\newcommand{\processHeight}{\textsf{processHeight}\xspace}
\newcommand{\processVote}{\textsf{processVote}\xspace}
\newcommand{\processInactivityPenalties}{\textsf{processInactivityPenalties}\xspace}
\newcommand{\finalizeConflict}{\textsf{finalizeConflict}\xspace}
\newcommand{\isSlashable}{\textsf{isSlashable}\xspace}
\newcommand{\constructVote}{\textsf{constructVote}\xspace}
\newcommand{\onBlock}{\textsf{onBlock}\xspace}
\newcommand{\updateForkChoiceRoot}{\textsf{updateForkChoiceRoot}\xspace}
\newcommand{\updateFinalized}{\textsf{updateFinalized}\xspace}
\newcommand{\getHead}{\textsf{getHead}\xspace}
\newcommand{\getConfirmed}{\textsf{getConfirmed}\xspace}
\newcommand{\getTarget}{\textsf{getTarget}\xspace}
\newcommand{\computeRoot}{\textsf{computeRoot}\xspace}

% Algorithmic event environment
\algdef{SE}[EVENT]{Event}{EndEvent}[1]{\textbf{at}\ #1\ \algorithmicdo}{\algorithmicend\ \textbf{event}}
\algtext*{EndEvent}

\pagestyle{plain}

\renewcommand{\paragraph}[1]{\vspace{7pt}\noindent\textbf{#1}}

\algrenewcommand\textproc{}

\usepackage{boxedminipage}

% Theorem environments
\theoremstyle{plain}
\newtheorem{theorem}{Theorem}
\newtheorem*{theorem*}{Theorem}
\newtheorem{corollary}{Corollary}
\newtheorem{lemma}{Lemma}
\newtheorem{proposition}{Proposition}
\newtheorem{property}[theorem]{Property}

\theoremstyle{definition}
\newtheorem{definition}[theorem]{Definition}

\theoremstyle{remark}
\newtheorem{remark}[theorem]{Remark}

\crefname{condition}{condition}{conditions}
\newtheorem{condition}{Condition}
\newtheorem{assumption}[theorem]{Assumption}
\newcommand{\GST}{\ensuremath{\mathsf{GST}}}


\title{Fresh Simplex with Timeouts}

\IfFileExists{lite_full_setup.tex}{% Shared lite/full setup for protocol-spec .tex files.
%
% Files under full/ that include this snippet compile in two modes:
%   - Full (default): all content rendered.
%   - Lite: when \litemode is defined before \documentclass (typically by a
%     wrapper .tex file at the repo root), lemmas, corollaries, propositions,
%     properties, proofs, remarks, assumptions, and conditions are dropped
%     via the comment package. Algorithms, theorem statements, definitions,
%     and surrounding prose remain.
%
% Place \input of this file in each spec's preamble AFTER all \newtheorem
% declarations (or just before \begin{document}). The exclusions use
% \@ifundefined so it's safe even when a file defines only a subset.
%
% Path resolution: \IfFileExists tries the local copy first (when cwd is
% full/), then falls back to full/lite_full_setup.tex (when cwd is the repo
% root, e.g. while compiling a top-level _lite.tex wrapper).
\newif\iffull
\ifx\litemode\undefined\fulltrue\else\fullfalse\fi
\iffull\else
  \usepackage{comment}
  \makeatletter
  \@ifundefined{lemma}{}{\excludecomment{lemma}}
  \@ifundefined{corollary}{}{\excludecomment{corollary}}
  \@ifundefined{proposition}{}{\excludecomment{proposition}}
  \@ifundefined{property}{}{\excludecomment{property}}
  \@ifundefined{proof}{}{\excludecomment{proof}}
  \@ifundefined{remark}{}{\excludecomment{remark}}
  \@ifundefined{assumption}{}{\excludecomment{assumption}}
  \@ifundefined{condition}{}{\excludecomment{condition}}
  \makeatother
\fi
}{% Shared lite/full setup for protocol-spec .tex files.
%
% Files under full/ that include this snippet compile in two modes:
%   - Full (default): all content rendered.
%   - Lite: when \litemode is defined before \documentclass (typically by a
%     wrapper .tex file at the repo root), lemmas, corollaries, propositions,
%     properties, proofs, remarks, assumptions, and conditions are dropped
%     via the comment package. Algorithms, theorem statements, definitions,
%     and surrounding prose remain.
%
% Place \input of this file in each spec's preamble AFTER all \newtheorem
% declarations (or just before \begin{document}). The exclusions use
% \@ifundefined so it's safe even when a file defines only a subset.
%
% Path resolution: \IfFileExists tries the local copy first (when cwd is
% full/), then falls back to full/lite_full_setup.tex (when cwd is the repo
% root, e.g. while compiling a top-level _lite.tex wrapper).
\newif\iffull
\ifx\litemode\undefined\fulltrue\else\fullfalse\fi
\iffull\else
  \usepackage{comment}
  \makeatletter
  \@ifundefined{lemma}{}{\excludecomment{lemma}}
  \@ifundefined{corollary}{}{\excludecomment{corollary}}
  \@ifundefined{proposition}{}{\excludecomment{proposition}}
  \@ifundefined{property}{}{\excludecomment{property}}
  \@ifundefined{proof}{}{\excludecomment{proof}}
  \@ifundefined{remark}{}{\excludecomment{remark}}
  \@ifundefined{assumption}{}{\excludecomment{assumption}}
  \@ifundefined{condition}{}{\excludecomment{condition}}
  \makeatother
\fi
}

\begin{document}
\maketitle

\section{Model and definitions}

We consider a system of~$n$ validators of which at most~$f$ are adversarial, where $n \geq 3f + 1$.
A \emph{quorum} is any set of at least $2n/3$ validators.
Any two quorums overlap in at least $n/3 + 1 > f$ validators, so their intersection contains at least one honest validator.

\begin{definition}[Height]\label{def:height}
\emph{Height} is the finality-gadget counter, separate from slot.
The genesis height is~$0$; height advances by exactly~$1$ via fixed-target justification (\Cref{def:justification}) or prefix notarization (\Cref{def:notarization}).
\end{definition}

\begin{definition}[Block]\label{def:block}
A \emph{block} has a \emph{parent} block, a \emph{slot} field, and a set of votes (\Cref{def:vote}).
The \emph{genesis block}~$\genesis$ has no parent and slot~$0$.
Slots strictly increase along parent links: $B.\slot > B.\parent.\slot$ for every non-genesis~$B$.
\end{definition}

\begin{definition}[Chain]\label{def:chain}
A \emph{chain} is a path from~$\genesis$ to some block~$B$ via parent links.
Chains are in bijection with their tips, and we identify the two.
Blocks form a tree rooted at~$\genesis$; we write $B \preceq B'$ when $B$ is an ancestor of~$B'$ (or $B = B'$), and $B \sim B'$ (\emph{compatible}) when $B \preceq B'$ or $B' \preceq B$.
Two chains \emph{conflict} when their tips are not compatible.
\end{definition}

\begin{definition}[Vote]\label{def:vote}
A \emph{vote} is a signed tuple
\[
  (\validator,\; \height,\; \slot,\; \target,\; \kind,\; \finalizeHeight,\; \finalizeTarget),
\]
where $\validator$ is the signer, $\height$ is a height, $\slot$ a slot, $\kind \in \{\just,\, \notar\}$, and $\target$ is a block, with $\target = \bot$ allowed only when $\kind = \notar$.
The pair $(\finalizeHeight,\, \finalizeTarget)$ is either both~$\bot$ or a commitment to finalize a block at a height.
A vote included on a chain must reference a target (and finalize target, when present) that exists on that chain, or be~$\bot$.
\end{definition}

\paragraph{State.}
The \emph{state after block~$B$}, written $\sigma[B]$, is the tuple
\[
  (L,\; s,\; h,\; s_h,\; \jtargets,\; \voteSlots,\; J,\; h_j,\; s_j,\; N,\; h_n,\; s_n,\; F,\; P),
\]
with components:
\begin{itemize}
  \item $L$: the \emph{chain head} ($L = B$ after processing~$B$).
  \item $s$: the \emph{current slot}, incremented once per slot by \processSlot (\Cref{alg:state-machine}).
  \item $h$: the \emph{current height}.
  \item $s_h$: the slot of the block whose processing last advanced~$h$ on this chain ($0$ if no advance has occurred).
  \item $\jtargets[1..n]$: for each validator~$i$, the target of~$i$'s \just vote at the current height ($\bot$ if none processed).
  \item $\voteSlots[1..n]$: for each~$i$, the maximum slot~$v.\slot$ over~$i$'s fresh votes at the current height, across both kinds ($\bot$ if none processed).
  \item $J$, $h_j$, $s_j$: the most recently justified block, its height, and its slot component (\Cref{def:justification}).
  \item $N$, $h_n$, $s_n$: the most recently notarized (or justified) block, its height, and its slot component; $J$ is itself a notarization (\Cref{def:notarization}).
  \item $F$: the most recently finalized block.
  \item $P \subseteq \{1, \ldots, n\}$: validators who have cast a vote with $\finalizeHeight = h_j$ and $\finalizeTarget = J$.
\end{itemize}
The genesis state is
\[
  \sigma_{\mathrm{gen}} \coloneqq (\genesis,\; 0,\; 0,\; 0,\; \bot^n,\; \bot^n,\; \genesis,\; 0,\; 0,\; \genesis,\; 0,\; 0,\; \genesis,\; \emptyset),
\]
with $\sigma[\genesis] \coloneqq \sigma_{\mathrm{gen}}$.
When unambiguous we drop the chain prefix and write $L$, $s$, $h$, etc.\ for the components of $\sigma[L]$.

\begin{definition}[State-height]\label{def:sheight}
The \emph{state-height} of a block~$B$ is $\sigma[B].h$: the height after processing~$B$ on its own chain.
\end{definition}

\paragraph{Height freshness.}
A \notar vote~$v$ is processed on a chain with head state $\sigma[L]$ only if $v.\height = \sigma[L].h$ and either $v.\target = \bot$ or $v.\target \preceq L$.
A \just vote~$v$ with target~$T$ is processed only if additionally
\[
  T \preceq L \;\wedge\; T.\slot \geq \sigma[L].s_h,
\]
so that $T$ lies on this chain and its slot is no earlier than where the chain's current height began.
Intuitively, a \just vote attests that its signer saw~$T$ as a chain of height~$v.\height$; a \notar vote attests only to the signer's current height, with the target signalling a commitment to~$T$ (if $T \neq \bot$) or no commitment (if $\bot$).

\begin{remark}[Freshness via $T$'s own chain]\label{rem:fresh-equiv}
For a \just vote, given $T \preceq L$, the constraint $T.\slot \geq \sigma[L].s_h$ is equivalent to $\sigma[T].h = \sigma[L].h$.
State-height is $\preceq$-monotone, equals $\sigma[L].h$ on~$L$, and advanced to its current value at the block whose slot is $\sigma[L].s_h$; so the ancestors of~$L$ at that state-height are exactly those with slot $\geq \sigma[L].s_h$.
\end{remark}

Freshness guarantees that every justification at height~$h$ is witnessed by \just votes whose targets' chains reached state-height~$h$, and that every notarization at height~$h$ is witnessed by votes cast at height~$h$.

Both justifications and notarizations carry a \emph{slot component}: the highest slot~$s$ for which a $2n/3$-quorum of contributing votes all have $\slot \geq s$.

\begin{definition}[Fixed-target justification]\label{def:justification}
A block~$T$ is \emph{justified at height~$h$ with slot~$s$} on a chain with state $\sigma[L]$ with $\sigma[L].h = h$ iff
\[
  \bigl|\{i \,:\, \jtargets[i] = T\}\bigr| \;\geq\; 2n/3
  \qquad\text{and}\qquad
  s \;=\; \max\bigl\{s' \;:\; |\{i : \jtargets[i] = T \,\wedge\, \voteSlots[i] \geq s'\}| \geq 2n/3\bigr\}.
\]
Only votes with target exactly~$T$ contribute.
\end{definition}

\begin{definition}[Notarization]\label{def:notarization}
A chain with state $\sigma[L]$ at height~$h$ is \emph{notarized at height~$h$ with slot~$s$} iff
\[
  \bigl|\{i \,:\, \voteSlots[i] \neq \bot\}\bigr| \;\geq\; 2n/3
  \qquad\text{and}\qquad
  s \;=\; \max\bigl\{s' \;:\; |\{i : \voteSlots[i] \geq s'\}| \geq 2n/3\bigr\}.
\]
Notarization does not require a common target; a \notar vote is fresh on this chain whenever $v.\height = h$ and $v.\target \in \{\bot\} \cup \{B : B \preceq L\}$. The notarized block is the chain head~$L$.
A justification of~$T$ is automatically a notarization of~$L$ at the same height: the $2n/3$ \just votes for~$T$ are fresh at height~$h$ on~$L$, so they contribute to $\voteSlots$; the slot component of the notarization coincides with that of the justification.
\end{definition}

\begin{definition}[Finality]\label{def:finality}
The justified block~$J$ is \emph{finalized at height~$h_j$} once $|P| \geq 2n/3$.
\end{definition}

\begin{definition}[Slashing, rule E1]\label{def:slashing}
A validator is \emph{slashable} if it signed two votes~$a, b$ (possibly $a = b$) with
\[
  b.\finalizeTarget \neq \bot \;\wedge\; a.\height = b.\finalizeHeight \;\wedge\; a.\target \neq b.\finalizeTarget.
\]
That is, a validator committing $(\finalizeHeight,\, \finalizeTarget) = (h_f, T_f)$ must not sign any vote at height~$h_f$ with target $\neq T_f$.
\end{definition}

Two \just (or two \notar) votes at the same height are not by themselves slashable; per-height uniqueness of honest \just votes is imposed behaviorally.
A validator who never sets $\finalizeTarget \neq \bot$ trivially avoids slashing.

\paragraph{State transitions.}
The state machine (\Cref{alg:state-machine}) exposes two per-chain routines: \processSlot, invoked once per slot on every maintained chain, and \processBlock, invoked additionally whenever a block~$B$ is processed.

\processSlot increments~$s$ and calls \processHeight, which in turn checks three branches in order: (i) finality, $F \gets J$ if $|P| \geq 2n/3$; (ii) justification, if some~$T$ has $\geq 2n/3$ entries in $\jtargets$; (iii) notarization, if $\geq 2n/3$ entries of $\voteSlots$ are non-$\bot$. Each height-advancing branch additionally records the slot component of the cert event.
Branches (ii) and (iii) both advance height and reset $\jtargets, \voteSlots$, so at most one height-advance fires per invocation.

$\processBlock(\sigma, B)$ has precondition $\sigma.s = B.\slot$: the caller runs \processSlot up to $B.\slot$ first.
Then $L \gets B$, and each vote~$v$ in~$B$ is fed through \processVote.
A \notar vote with $v.\height = \sigma.h$ and $v.\target \in \{\bot\} \cup \{B : B \preceq \sigma.L\}$ updates $\voteSlots[i] \gets \max(\voteSlots[i],\, v.\slot)$ (treating~$\bot$ as $-\infty$); a \just vote additionally passes the target-freshness check $v.\target \preceq \sigma.L \wedge v.\target.\slot \geq \sigma.s_h$ and, if it does, updates both $\voteSlots[i]$ and $\jtargets[i] \gets v.\target$.
The $P$-gate adds~$i$ to~$P$ iff $v.\finalizeHeight = \sigma.h_j$ and $v.\finalizeTarget = \sigma.J$; this is independent of freshness.
\processBlock never advances~$h$: height advances happen only in \processSlot, which runs on every slot (including empty ones).

\begin{figure}[H]
\caption{State machine}\label{alg:state-machine}
\begin{center}
\begin{boxedminipage}{\textwidth}
\begin{algorithmic}[1]

\Function{\processSlot}{$\sigma$} \Comment{Runs once per slot on every chain, even without a block}
  \State $\sigma.s \gets \sigma.s + 1$
  \State $\sigma \gets \processHeight(\sigma)$
  \State \Return $\sigma$
\EndFunction
\vspace{3pt}

\Function{\processBlock}{$\sigma,\; B$}
  \State \textbf{assert} $\sigma.s = B.\slot$ \Comment{Precondition: caller advanced $\sigma.s$ via \processSlot}
  \State $\sigma.L \gets B$
  \For{each vote $v$ in $B$} $\sigma \gets \processVote(\sigma,\, v)$ \EndFor
  \State \Return $\sigma$
\EndFunction
\vspace{3pt}

\Function{\processVote}{$\sigma,\; v$}
  \State $i \gets v.\validator$
  \If{$v.\height = \sigma.h$ \textbf{and} $(v.\target = \bot$ \textbf{or} $v.\target \preceq \sigma.L)$} \Comment{enough for \notar}
    \State $\sigma.\voteSlots[i] \gets \max(\sigma.\voteSlots[i],\, v.\slot)$
    \If{$v.\kind = \just$ \textbf{and} $v.\target \neq \bot$ \textbf{and} $v.\target.\slot \geq \sigma.s_h$}
      \State $\sigma.\jtargets[i] \gets v.\target$ \Comment{\just additionally requires target freshness}
    \EndIf
  \EndIf
  \If{$v.\finalizeHeight = \sigma.h_j$ \textbf{and} $v.\finalizeTarget = \sigma.J$}
    \State $\sigma.P \gets \sigma.P \cup \{i\}$
  \EndIf
  \State \Return $\sigma$
\EndFunction
\vspace{3pt}

\Function{\processHeight}{$\sigma$}
  \If{$|\sigma.P| \geq 2n/3$} \Comment{Finality}
    \State $\sigma.F \gets \sigma.J$
  \EndIf
  \State $T \gets \arg\max_{T' \in \{\sigma.\jtargets[i] \,:\, \sigma.\jtargets[i] \neq \bot\}} |\{i : \sigma.\jtargets[i] = T'\}|$ \Comment{$T \gets \bot$ if the set is empty;}
  \If{$T \neq \bot$ \textbf{and} $|\{i : \sigma.\jtargets[i] = T\}| \geq 2n/3$} \Comment{Justification (also a notarization)}
    \State $s^\star \gets \max\{s' : |\{i : \sigma.\jtargets[i] = T \land \sigma.\voteSlots[i] \geq s'\}| \geq 2n/3\}$
    \State $\sigma.J \gets T$, \; $\sigma.h_j \gets \sigma.h$, \; $\sigma.s_j \gets s^\star$, \; $\sigma.P \gets \emptyset$
    \State $\sigma.N \gets \sigma.L$, \; $\sigma.h_n \gets \sigma.h$, \; $\sigma.s_n \gets s^\star$
    \State $\sigma.h \gets \sigma.h + 1$, \; $\sigma.s_h \gets \sigma.L.\slot$
    \State reset $\jtargets, \voteSlots$
    \State \Return $\sigma$
  \EndIf
  \If{$|\{i : \sigma.\voteSlots[i] \neq \bot\}| \geq 2n/3$} \Comment{Notarization (no target check)}
    \State $s^\star \gets \max\{s' : |\{i : \sigma.\voteSlots[i] \geq s'\}| \geq 2n/3\}$
    \State $\sigma.N \gets \sigma.L$, \; $\sigma.h_n \gets \sigma.h$, \; $\sigma.s_n \gets s^\star$
    \State $\sigma.h \gets \sigma.h + 1$, \; $\sigma.s_h \gets \sigma.L.\slot$
    \State reset $\jtargets, \voteSlots$
    \State \Return $\sigma$
  \EndIf
  \State \Return $\sigma$
\EndFunction
\vspace{3pt}

\Function{\finalizeConflict}{$v_1, v_2$} \Comment{Does $v_1$ conflict with $v_2$'s finalize commitment?}
  \State \Return $v_2.\finalizeTarget \neq \bot \;\wedge\; v_1.\height = v_2.\finalizeHeight \;\wedge\; v_1.\target \neq v_2.\finalizeTarget$
\EndFunction
\vspace{3pt}

\Function{\isSlashable}{$v_1, v_2$}
  \State \Return $v_1.\validator = v_2.\validator$
  \Statex \hspace{3em}$\land \;(\finalizeConflict(v_1, v_2) \lor \finalizeConflict(v_2, v_1))$
\EndFunction

\end{algorithmic}
\end{boxedminipage}
\end{center}
\end{figure}

\section{Accountable safety}

\begin{lemma}[Height progression]\label{lem:heightprog}
\processHeight increments~$h$ by at most~$1$ per invocation, and each increment is gated by a $\geq 2n/3$ justification or notarization quorum.
\end{lemma}

\begin{proof}
The only assignments to~$h$ are $h \gets h + 1$ in the justification and notarization branches of \processHeight, each guarded by its $\geq 2n/3$ quorum.
\end{proof}

\begin{lemma}[Checkpoint monotonicity]\label{lem:slotmono}
On any chain, the fields $s$, $h$, $h_j$, $h_n$, $s_h$, $F$, $J$, $N$ are non-decreasing along~$\preceq$, with $F \preceq J \preceq N$ at all times.
$J$ and $h_j$ advance only on justification events, and $h_n, s_h$ only on height-advances, with $s_h$ set to the slot of the advancing block.
Consequently, when $\sigma[Y].h = h$, the ancestors of~$Y$ at state-height~$h$ are $\{B' \preceq Y : B'.\slot \geq \sigma[Y].s_h\}$.
\end{lemma}

\begin{proof}
Reading off \Cref{alg:state-machine}: $s$ only increases (in \processSlot); $h$ only increases, by~$1$, in the height-advance branches of \processHeight, which also set $h_j, h_n \gets h$ on their respective updates and $s_h \gets L.\slot$.
Slots strictly increase along~$\preceq$ (\Cref{def:block}), so $s_h$ is non-decreasing.
$F$ updates only via $F \gets J$, so $F \preceq J$ reduces to $J$-monotonicity.

\emph{Height-$h$ interval.}
Along any chain, $h$ increases monotonically; when $\sigma[Y].h = h$, the value $s_h$ was set at the block that advanced state-height to~$h$.
By strict slot increase along~$\preceq$ and the monotonicity of $(h, s_h)$, the ancestors $B' \preceq Y$ with $\sigma[B'].h = h$ are exactly those with $B'.\slot \geq s_h$.

\emph{Block fields $J$, $N$.}
Consider a height-advance on chain~$Y$ with pre-state height~$h$.
If the justification branch fires, every contributing \just vote has target~$T$ and passes target freshness, forcing $T \preceq L$ and $T.\slot \geq s_h$; the height-$h$ interval claim gives $\sigma[T].h = h$.
If the notarization branch fires, no target constraint is imposed, but $\sigma[L].h = h$ directly.
In both cases the previous notarized block $N_{\mathrm{old}}$ satisfies $\sigma[N_{\mathrm{old}}].h \leq h - 1$ and $N_{\mathrm{old}} \preceq Y$, so $N_{\mathrm{old}} \prec L$ (and $N_{\mathrm{old}} \prec T \preceq L$ in the justification case) by state-height monotonicity.
The justification branch sets $J \gets T$ and $N \gets L$ with $T \preceq L$; the notarization branch sets only $N \gets L$ and leaves $J$ unchanged.
Hence $N$ strictly advances at every height-advance, $J$ at every justification, and $J \preceq N$ is preserved.
\end{proof}

\begin{lemma}[Any chain past a finalized height contains the finalized block]\label{lem:mainsafety}
Unless $\geq n/3$ validators are slashable: if~$C$ is finalized at height~$h$, then $C \preceq B$ for every block~$B$ with $\sigma[B].h > h$.
\end{lemma}

\begin{proof}
By \Cref{lem:heightprog}, some $B^\star \preceq B$ advanced height from~$h$ to~$h+1$ via a quorum~$Q$ of votes at height~$h$ processed on~$B^\star$'s chain; each $i \in Q$ has at least one such vote~$a_i$ included in some block $B_{a_i} \preceq B^\star$.
Finality of~$C$ at~$h$ gives a quorum~$Q_F$ each signing a vote~$b$ with $b.\finalizeHeight = h$ and $b.\finalizeTarget = C$.

By quorum intersection, $|Q \cap Q_F| \geq n/3$; pick $v \in Q \cap Q_F$ not slashable.
E1 (\Cref{def:slashing}) applied to $v$'s votes $a \in Q$, $b \in Q_F$ with $a.\height = h = b.\finalizeHeight$ forces $a.\target = C$.
By \Cref{def:vote}, $a.\target$ lies on the chain containing $B_a$, so $C = a.\target \preceq B_a \preceq B^\star \preceq B$.
\end{proof}

\begin{lemma}[Finalized blocks form a chain]\label{lem:finchain}
Unless $\geq n/3$ validators are slashable: any two finalized checkpoints $(C, h), (C', h')$ are compatible, with $C \preceq C'$ whenever $h \leq h'$.
\end{lemma}

\begin{proof}
\emph{Case $h = h'$.}
Finality of~$C$ gives a quorum~$Q_F$ with $\finalizeHeight = h$, $\finalizeTarget = C$.
Finality of~$C'$ requires justification of~$C'$, giving a \just-quorum~$Q'$ at height~$h$ with target~$C'$.
By quorum intersection, $|Q_F \cap Q'| \geq n/3$; pick $v \in Q_F \cap Q'$ not slashable.
E1 (\Cref{def:slashing}) applied to $v$'s votes in~$Q_F, Q'$ forces $C = C'$.

\emph{Case $h < h'$.}
Let~$Z$ be the chain on which~$C'$ was finalized, at the post-state where the finalize-quorum closed; then $\sigma[Z].h \geq h' > h$, $C' \preceq Z$, and $\sigma[C'].h = h'$ (\Cref{rem:fresh-equiv}).
\Cref{lem:mainsafety} at~$Z$ gives $C \preceq Z$; \Cref{rem:fresh-equiv} on~$C$'s own finalizing chain gives $\sigma[C].h = h$.
Both~$C, C'$ lie on~$Z$, so comparable; with $\sigma[C].h = h < h' = \sigma[C'].h$, state-height monotonicity (\Cref{lem:slotmono}) forces $C \prec C'$.
\end{proof}

\begin{theorem}[Accountable safety]\label{thm:safety}
Unless $\geq n/3$ validators are slashable, no two conflicting blocks can be finalized.
\end{theorem}

\begin{proof}
Immediate from \Cref{lem:finchain}: any two finalized blocks are compatible.
\end{proof}

\section{Fork-choice store}\label{sec:store}

A node may track multiple chains simultaneously; the \emph{store} is the node-local data structure that mediates which chain the node follows.
It maintains a single \emph{root block} $R$ alongside a tracker $J$ of the highest justified block ever observed (extending the current store-finalized block).
On each block acceptance, $R$ is recomputed from $J$ and the block tree by the following rule, where $\Sigma.h^\star$ is the highest state-height over the leaves of $\Sigma.\T$ (maintained incrementally by \onBlock).
\begin{itemize}
  \item If $\sigma[J].h \geq \Sigma.h^\star - 1$ and $\Sigma.F \preceq J$, set $R \gets J$.
  \item Otherwise, set $R$ to the block in $\Sigma.\T$ at state-height $\Sigma.h^\star - 1$ extending $\Sigma.F$ with the largest $\hash$ (ties on hash do not occur); if that set is empty, leave $R$ unchanged.
\end{itemize}

\begin{definition}[Cert event]\label{def:cert-event}
When \processHeight fires a height advance on a chain (\Cref{alg:state-machine}), we call the triggering event a \emph{cert event}, characterized by:
\begin{itemize}
  \item a \emph{root block} $X$: the justified target~$T$ in the justification branch, and the chain head~$L$ at firing time in the notarization branch;
  \item the height~$h$ at which the advance fires, equal to the pre-advance $\sigma.h$ and to the post-advance $\sigma.h_n$;
  \item a \emph{bit} $b \in \{0, 1\}$: $b = 1$ for a justification, $b = 0$ for a notarization without accompanying justification.
\end{itemize}
\end{definition}

\begin{definition}[Store]\label{def:store}
A node maintains a \emph{store}
\[
  \Sigma \;=\; (\sigma,\; \T,\; F,\; R,\; J,\; h^\star),
\]
where:
\begin{itemize}
  \item $\sigma$ is the \emph{state map}, assigning each accepted block its per-chain post-state (\Cref{alg:state-machine});
  \item $\T$ is the \emph{block tree}, the set of accepted blocks;
  \item $F$ is the \emph{store-finalized block};
  \item $R$ is the \emph{store root};
  \item $J$ is the \emph{highest justified block}: a block, extending $\Sigma.F$, with the largest state-height $\sigma[J].h$ among justifications observed on any processed chain;
  \item $h^\star$ is the highest state-height over the leaves of $\Sigma.\T$, maintained incrementally by \onBlock.
\end{itemize}
The genesis store has $\T = \{\genesis\}$, $F = R = J = \genesis$, $h^\star = 0$, treating genesis as a justification at height~$0$.
\end{definition}

The capital~$\Sigma$ distinguishes the store from the per-chain state~$\sigma$.
We write $\Sigma.\sigma$, $\Sigma.\T$, $\Sigma.F$, $\Sigma.R$, $\Sigma.J$, $\Sigma.h^\star$ for field access.

The acceptance routine \onBlock (\Cref{alg:store}) admits $B$ only if $B.\parent \in \Sigma.\T$ and $\Sigma.F \preceq B$.
It runs \processSlot once per slot from $\Sigma.\sigma[B.\parent].s + 1$ up to $B.\slot$ (so any height advance enabled by the accumulated tally fires before~$B$'s votes are folded in), then $\processBlock(\cdot, B)$; the resulting post-state becomes $\Sigma.\sigma[B]$.
The post-state then offers the justified-checkpoint descriptor $(\sigma[B].h_j,\, \sigma[B].J)$ to \updateForkChoiceRoot, which (i) raises $\Sigma.J$ to the new justified block if it has a strictly higher state-height than the current $\Sigma.J$ and extends $\Sigma.F$, and (ii) recomputes $\Sigma.R$ by the rule above.
Finally, \updateFinalized sets $\Sigma.F \gets \sigma[B].F$ iff $\sigma[B].F \succ \Sigma.F$ and $\sigma[B].F \preceq \Sigma.R$.

$\getConfirmed(\Sigma, \Omega)$ returns a block on the available chain that, under synchrony, is safe in the sense that it will not be reverted as long as more than half of the online validators are honest. Here $\Omega$ represents the state of the available chain --- whatever extra information the validator uses to disambiguate among descendants of $\Sigma.R$ in $\Sigma.\T$. The protocol fixes only that the returned block lies in $\Sigma.\T$ and descends from $\Sigma.R$; it does not constrain $\Omega$'s structure or how the choice depends on it. Different validators may pass different~$\Omega$, and a single validator may pass different~$\Omega$ at different call sites.

\begin{figure}[H]
\caption{Store and fork-choice root}\label{alg:store}
\begin{center}
\begin{boxedminipage}{\textwidth}
\begin{algorithmic}[1]

\State \textbf{Genesis store} $\Sigma$:\ $\Sigma.\T \gets \{\genesis\},\ \Sigma.F \gets \genesis,\ \Sigma.R \gets \genesis,\ \Sigma.J \gets \genesis,\ \Sigma.h^\star \gets 0,\ \Sigma.\sigma[\genesis] \gets \sigma_{\mathrm{gen}}$
\vspace{3pt}

\Function{\onBlock}{$\Sigma,\; B$}
  \State \textbf{assert} $B.\parent \in \Sigma.\T$
  \State \textbf{assert} $\Sigma.F \preceq B$
  \State $\Sigma.\T \gets \Sigma.\T \cup \{B\}$
  \State $\sigma' \gets \Sigma.\sigma[B.\parent]$
  \While{$\sigma'.s < B.\slot$} \Comment{Advance through every intervening slot}
    \State $\sigma' \gets \processSlot(\sigma')$
  \EndWhile
  \State $\sigma' \gets \processBlock(\sigma',\; B)$
  \State $\Sigma.\sigma[B] \gets \sigma'$
  \State $\Sigma.h^\star \gets \max(\Sigma.h^\star,\; \sigma'.h)$ \Comment{$\sigma'.h$ is $B$'s state-height}
  \State $\Sigma \gets \updateForkChoiceRoot(\Sigma,\; (\sigma'.h_j,\; \sigma'.J))$
  \State $\Sigma \gets \updateFinalized(\Sigma,\; \sigma'.F)$
  \State \Return $\Sigma$
\EndFunction
\vspace{3pt}

\Function{\updateForkChoiceRoot}{$\Sigma,\; (h,\, J)$}
  \If{$h > \sigma[\Sigma.J].h$ \textbf{and} $\Sigma.F \preceq J$} \Comment{Raise $J$ when dominated and on the $F$-chain}
    \State $\Sigma.J \gets J$
  \EndIf
  \If{$\sigma[\Sigma.J].h \geq \Sigma.h^\star - 1$ \textbf{and} $\Sigma.F \preceq \Sigma.J$}
    \State $\Sigma.R \gets \Sigma.J$
  \Else
    \State $\Sigma.R \gets \arg\max_{B \in \Sigma.\T,\; \sigma[B].h = \Sigma.h^\star - 1,\; B \succeq \Sigma.F} \hash(B)$ \Comment{empty $\Rightarrow$ $\Sigma.R$ unchanged}
  \EndIf
  \State \Return $\Sigma$
\EndFunction
\vspace{3pt}

\Function{\updateFinalized}{$\Sigma,\; F_B$}
  \If{$F_B \succ \Sigma.F$ and $F_B \preceq \Sigma.R$}
    \State $\Sigma.F \gets F_B$
  \EndIf
  \State \Return $\Sigma$
\EndFunction
\vspace{3pt}

\Function{\getConfirmed}{$\Sigma,\; \Omega$} \Comment{Choice depends on auxiliary state~$\Omega$}
  \State \Return a block in~$\Sigma.\T$ that descends from~$\Sigma.R$, depending on~$\Omega$
\EndFunction

\end{algorithmic}
\end{boxedminipage}
\end{center}
\end{figure}

The $F$-filter never discards a justified block that should reach $J$: \Cref{thm:fleqr} below shows that $F$ advances only to blocks $\preceq R$, so every justified block that could beat the current $J$ on height remains on the $F$-chain.

\subsection{Store invariants and safety}

\subsubsection*{Unconditional invariants}

\begin{lemma}[$J$ height monotonicity]\label{lem:J-monotone}
The state-height $\sigma[\Sigma.J].h$ is non-decreasing over time.
\end{lemma}

\begin{proof}
$\Sigma.J$ changes only inside \updateForkChoiceRoot, whose guard $h > \sigma[\Sigma.J].h$ allows only strict increases of the height; $\sigma[\cdot].h$ on a fixed block is invariant once set.
\end{proof}

\begin{theorem}[Local irreversibility of finality]\label{thm:finperm}
Once a node sets $F = F$, $F$ descends from~$F$ at all future times.
\end{theorem}

\begin{proof}
$F$ is updated only by \updateFinalized, whose guard $F_B \succ \Sigma.F$ forces strict extension along~$\preceq$.
\end{proof}

\begin{theorem}[$F \preceq R$]\label{thm:fleqr}
The store maintains $F \preceq R$ at all times.
\end{theorem}

\begin{proof}
By induction on store operations.
At genesis, $F = R = \genesis$, so $F \preceq R$.
Assume the invariant before a call to either mutator in \Cref{alg:store}.

\emph{\updateForkChoiceRoot.}
$F$ is untouched. The recompute sets $R$ to either $\Sigma.J$ (only when $\Sigma.F \preceq \Sigma.J$, by the explicit branch guard) or to a block at state-height $h^\star - 1$ filtered by $\succeq \Sigma.F$; if neither branch fires (the fallback set being empty), $R$ stays unchanged. In all three cases $R \succeq F$ holds, the third by induction.

\emph{\updateFinalized.}
$R$ is untouched, and the guard $F_B \preceq \Sigma.R$ ensures $F \gets F_B \preceq R$.
\end{proof}

\begin{remark}[Store-finalized invariant]\label{rem:fs-invariant}
At all times $\Sigma.F$ is either~$\genesis$ or a block that was finalized on some chain; moreover, whenever a justification $(J, h)$ has fired on some processed chain, the descriptor $(h,\, J)$ was offered to \updateForkChoiceRoot, so $\sigma[\Sigma.J].h \geq h$ for every such $(J, h)$ with $J \succeq \Sigma.F$ (\Cref{lem:J-monotone}).
$\Sigma.F$ only advances via \updateFinalized to $F_B = \sigma[B].F$, and $\sigma[B].F \neq \genesis$ requires the finality branch of \processHeight to have fired on~$B$'s chain, so $F_B$ is a block finalized on that chain.
\end{remark}

\begin{theorem}[Fork-choice consistency]\label{thm:fcconsistency}
Once a node sets $F = F$, $\getConfirmed(\Sigma, \Omega)$ returns a block descending from~$F$ at all future times, for every~$\Omega$.
\end{theorem}

\begin{proof}
By \Cref{thm:finperm} and~\Cref{thm:fleqr}, $F \preceq F \preceq R$, and $\getConfirmed(\Sigma, \Omega)$ returns a descendant of~$R$ regardless of~$\Omega$.
\end{proof}

\subsubsection*{Further properties under $f < n/3$}

\begin{lemma}[Checkpoint chain]\label{lem:certchain}
Unless $\geq n/3$ validators are slashable: if~$F$ is finalized at height~$h_F$ and a cert event $(C,\, h,\, s,\, b)$ has fired on any chain processed by the node with $h \geq h_F$, then $F$ and~$C$ are compatible, with $F \prec C$ when $h > h_F$.
\end{lemma}

\begin{proof}
If $h_F = 0$ then $F = \genesis$ and compatibility is trivial; assume $h_F \geq 1$.
The cert event fired at the post-state of some processed block~$B$ whose \processSlot call triggered \processHeight with $(\sigma[B].N,\, \sigma[B].h_n) = (C,\, h)$: for the justification branch $C = T \preceq L$ via target freshness, and for the notarization branch $C = L$ directly; in both cases $C \preceq B$.
Every $h_n$-update is paired with $\sigma.h \gets \sigma.h + 1$ (\Cref{lem:slotmono}), hence $\sigma[B].h \geq h + 1 > h_F$; \Cref{lem:mainsafety} gives $F \preceq B$, and with $C \preceq B$ this yields $F \sim C$.
For $h > h_F$: $\sigma[F].h = h_F < h = \sigma[C].h$ (justification: target freshness; notarization: pre-advance $\sigma[C].h = h$), and state-height monotonicity along~$\preceq$ (\Cref{lem:slotmono}) forces $F \prec C$ on the comparable pair $F \sim C$.
\end{proof}

\begin{lemma}[Upgrade property]\label{lem:upgrade}
Unless $\geq n/3$ validators are slashable: if~$F$ is finalized at height~$h_F$ and some block~$B$ whose post-state has $(F, h_F)$ as its justified checkpoint has been processed by the node, then $R \succeq F$ at all future times.
\end{lemma}

\begin{proof}
Fix the moment when $\onBlock(B)$ offered the descriptor $(h_F,\, F)$ to \updateForkChoiceRoot, and let $F^0 = \Sigma.F$ at that moment.
Both $F^0$ and $F$ are finalized or genesis (\Cref{rem:fs-invariant}), so \Cref{lem:finchain} gives $F^0 \preceq F$ or $F \preceq F^0$.

\emph{Case $F \preceq F^0$.}
$R \succeq F^0 \succeq F$ already holds by \Cref{thm:fleqr}; \Cref{thm:finperm} preserves $F \succeq F^0$ and \Cref{thm:fleqr} preserves $R \succeq F$ at all future times, so $R \succeq F$ persists.

\emph{Case $F^0 \preceq F$.}
The filter $\Sigma.F \preceq F$ passes, so the call sets $\Sigma.J$ to a block of state-height $\geq h_F$; \Cref{lem:J-monotone} preserves $\sigma[\Sigma.J].h \geq h_F$ thereafter.
Write $h_J \coloneqq \sigma[\Sigma.J].h$.
If $h_J > h_F$, \Cref{lem:certchain} applied to the justification cert event at $\Sigma.J$ gives $F \prec \Sigma.J$.
If $h_J = h_F$, $\Sigma.J$ was justified at height~$h_F$ on some processed chain, giving a \just-quorum~$Q'$ of size $\geq 2n/3$ targeting $\Sigma.J$ at height~$h_F$; finality of~$F$ at~$h_F$ gives a quorum~$Q_F$ with $\finalizeHeight = h_F$ and $\finalizeTarget = F$, and quorum intersection plus E1 (as in the proof of \Cref{lem:finchain}) forces $\Sigma.J = F$.
Either way $\Sigma.J \succeq F$, hence $\sigma[\Sigma.J].h \geq h_F$ and $\Sigma.F \preceq F \preceq \Sigma.J$.

It remains to argue that the \updateForkChoiceRoot recompute in any future call yields $R \succeq F$.
Let $h^\star$ be the current leaf-height bound, and consider the two branches.
If the branch $R \gets \Sigma.J$ fires, $R = \Sigma.J \succeq F$.
Otherwise the fallback fires: $R$ is set to a block at state-height $h^\star - 1$ extending $\Sigma.F$ (or stays unchanged if no such block exists).
Note $\sigma[F].h = h_F$ and $h^\star \geq h_F + 1$ (since $\Sigma.\T$ contains a block witnessing $F$'s finalization at state-height $\geq h_F + 1$), so $h^\star - 1 \geq h_F$; \Cref{lem:mainsafety} (when $h^\star - 1 > h_F$) and direct identity (when $h^\star - 1 = h_F$, in which case $F$ itself is a candidate) give $R \succeq F$.
If the fallback set is empty, $R$ is unchanged from its previous value, which by induction was $\succeq F$.
\end{proof}

\begin{theorem}[Local acceptance of finality updates]\label{thm:finlive}
Unless $\geq n/3$ validators are slashable: if a block~$B$ is processed by \onBlock and its post-state has finalized block~$F_B$, then after processing $F \succeq F_B$.
\end{theorem}

\begin{proof}
If $F \succeq F_B$ already, the claim holds.
Otherwise $F$ is genesis or finalized (\Cref{rem:fs-invariant}), so \Cref{lem:finchain} gives $F \sim F_B$ and hence $F \prec F_B$.
$F_B$ being finalized in~$B$'s post-state requires some ancestor~$B' \preceq B$ to have post-state $(J_{B'},\, h_j^{B'}) = (F_B,\, h_{F_B})$; $B'$ was processed before~$B$ (parent-first), so \Cref{lem:upgrade} gives $R \succeq F_B$, and the \updateFinalized guard passes.
\end{proof}

\begin{theorem}[Lock-in]\label{thm:lockin}
Unless $\geq n/3$ validators are slashable: if block~$F$ is finalized at height~$h$ on any chain, and some block whose post-state has $(F, h)$ as its justified checkpoint has been processed by the node, then $R \succeq F$ at all future times, and $\getConfirmed(\Sigma, \Omega)$ always returns a descendant of~$F$, for every~$\Omega$.
\end{theorem}

\begin{proof}
\Cref{lem:upgrade} gives $R \succeq F$ directly.
$\getConfirmed(\Sigma, \Omega)$ (descending from~$R$ for every~$\Omega$) therefore returns a descendant of~$F$.
\end{proof}

\begin{theorem}[Order independence]\label{thm:orderindep}
Unless $\geq n/3$ validators are slashable: the store state after processing a set of blocks depends only on which blocks have been processed, not on their order.
\end{theorem}

\begin{proof}
Each block's post-state is a deterministic function of itself and its parent's post-state, so $\sigma$ depends only on the processed set.

\emph{$F$.} By \Cref{lem:finchain}, $\{\sigma[B].F : B \in \T\}$ is a chain with maximum $F_{\max}$, attained at some processed~$B^\ast$.
\Cref{rem:fs-invariant} forces $F$ into this chain, so $F \preceq F_{\max}$; once $B^\ast$ has been processed, \Cref{thm:finlive} gives $F \succeq F_{\max}$.

\emph{$\Sigma.J$.} The multiset~$D$ of justified-checkpoint descriptors $(h, J)$ offered to \updateForkChoiceRoot is a deterministic function of the processed set: each block contributes one such descriptor exactly once.
By \Cref{lem:just-unique-height}, each height has at most one justified block on any processed chain, so the height-to-block map $h \mapsto J$ over $D$ is well-defined; \Cref{lem:certchain} together with the $\Sigma.F \preceq J$ filter ensures every descriptor admitted into the running max extends $F_{\max}$.
The final $\Sigma.J$ is therefore the unique block of largest height in $D$ extending $F_{\max}$, a deterministic function of the processed set.

\emph{$R$.} The recompute in \updateForkChoiceRoot is a pure function of $(\Sigma.\T, \sigma, \Sigma.F, \Sigma.J)$: the leaf set, $h^\star$, the $\sigma[B].h$ values, and the $\succeq F$-filtered candidate set at $h^\star - 1$ all depend only on the processed set; the hash-argmax is deterministic. Hence $R$'s final value is order-independent.
\end{proof}

\begin{corollary}[Cross-node agreement]\label{cor:agreement}
Unless $\geq n/3$ validators are slashable: two nodes that have processed the same set of blocks have the same $(\sigma,\, \T,\, F,\, R,\, J)$.
\end{corollary}

\begin{proof}
Immediate from \Cref{thm:orderindep}.
\end{proof}

\section{Inactivity leak}

This section extends the per-chain state machine of \Cref{alg:state-machine} with a penalty counter and a new transition function \processInactivityPenalties, invoked each slot by an extended \processSlot (\Cref{alg:leak-processslot}); \processBlock, \processHeight, and \processVote are unchanged.
Write $X = \getConfirmed(\Sigma, \Omega)$ for the \emph{canonical chain}, with $\Omega$ the validator's auxiliary state (\Cref{alg:store}).

Two results follow.
\emph{Tightness} (\Cref{thm:tightness}) is purely state-logic: on any fixed chain, the total penalty increment over a non-finality period is at least $(T-1) \cdot n/3$, with no hypothesis on synchrony, honesty, or the slashable set.
\emph{Fairness} (\Cref{thm:fairness}) is local to a single honest~$i$: the vote~$i$ produces per the honest rule (\Cref{def:voting-rule}), once included on any extension of~$i$'s canonical chain still at the same state-height, exempts~$i$ from Layer~1 penalties at that state-height.
Layer-2 exemption is out of scope: it depends on quorum formation on a common target, not on any local property of~$i$.

\subsection{Penalty units}

Augment the per-chain state tuple with one non-negative counter $\sigma.\penalties[1..n]$, initialized to~$0$ at genesis; no other field is added or modified.
\Cref{alg:leak-processslot} replaces \processSlot from \Cref{alg:state-machine} with a version that snapshots the pre-advance state and invokes \processInactivityPenalties after \processHeight.
The store of \Cref{sec:store} invokes this extended \processSlot in place of the base one.

\begin{figure}[H]
\caption{Penalty-tracking state extension (leak section)}\label{alg:leak-processslot}
\begin{center}
\begin{boxedminipage}{\textwidth}
\begin{algorithmic}[1]

\Function{\processSlot}{$\sigma$} \Comment{extended: tracks penalties}
  \State $\sigma.s \gets \sigma.s + 1$
  \State $\sigma_{\mathrm{pre}} \gets \sigma$ \Comment{snapshot of pre-advance state}
  \State $\sigma \gets \processHeight(\sigma)$
  \State $\sigma \gets \processInactivityPenalties(\sigma_{\mathrm{pre}},\; \sigma)$
  \State \Return $\sigma$
\EndFunction
\vspace{3pt}

\Function{\processInactivityPenalties}{$\sigma_{\mathrm{pre}},\; \sigma_{\mathrm{post}}$}
  \If{$\sigma_{\mathrm{pre}}.h = \sigma_{\mathrm{post}}.h$} \Comment{Layer 1: stall}
    \For{each $i \in \{1, \ldots, n\}$}
      \If{$\sigma_{\mathrm{post}}.\voteSlots[i] = \bot$}
        \State $\sigma_{\mathrm{post}}.\penalties[i] \gets \sigma_{\mathrm{post}}.\penalties[i] + 1$
      \EndIf
    \EndFor
  \EndIf
  \If{$\sigma_{\mathrm{pre}}.J = \sigma_{\mathrm{post}}.J$} \Comment{Layer 2 target: justification did \emph{not} advance}
    \State $T^* \gets \arg\max_{T' \in \{\sigma_{\mathrm{pre}}.\jtargets[i] \,:\, \sigma_{\mathrm{pre}}.\jtargets[i] \neq \bot\}} |\{i : \sigma_{\mathrm{pre}}.\jtargets[i] = T'\}|$ \Comment{$T^* \gets \bot$ if the set is empty; ties broken deterministically}
    \For{each $i \in \{1, \ldots, n\}$}
      \If{$\sigma_{\mathrm{pre}}.\jtargets[i] \neq T^*$ \textbf{or} $\sigma_{\mathrm{pre}}.\jtargets[i] = \bot$}
        \State $\sigma_{\mathrm{post}}.\penalties[i] \gets \sigma_{\mathrm{post}}.\penalties[i] + 1$
      \EndIf
    \EndFor
  \EndIf
  \If{$\sigma_{\mathrm{pre}}.F = \sigma_{\mathrm{post}}.F$ \textbf{and} $\sigma_{\mathrm{pre}}.F \prec \sigma_{\mathrm{pre}}.J$} \Comment{Layer 2 finalize: $F$ did not advance, pending finality}
    \For{each $i \in \{1, \ldots, n\}$}
      \If{$i \notin \sigma_{\mathrm{pre}}.P$}
        \State $\sigma_{\mathrm{post}}.\penalties[i] \gets \sigma_{\mathrm{post}}.\penalties[i] + 1$
      \EndIf
    \EndFor
  \EndIf
  \State \Return $\sigma_{\mathrm{post}}$
\EndFunction

\end{algorithmic}
\end{boxedminipage}
\end{center}
\end{figure}

\paragraph{Reading \Cref{alg:leak-processslot}.}
\processInactivityPenalties checks three independent guards; any subset may fire at a given slot.

\emph{Layer~1 (stall).} Fires when the slot did not advance height.
Every~$i$ with $\sigma_{\mathrm{post}}.\voteSlots[i] = \bot$ is bumped, i.e.\ every validator that failed to have any fresh vote registered at the current height ($\voteSlots$ is not reset at a stall).

\emph{Layer~2 target.} Fires when the justification branch of \processHeight did not fire (else $J$ would strictly advance by \Cref{lem:slotmono}).
$T^*$ is the plurality target across $\sigma_{\mathrm{pre}}.\jtargets$ (deterministic tiebreak), or $\bot$ if no \just votes have been registered; the second disjunct of the guard bumps non-voters when $T^* = \bot$ (in which case all~$n$ validators are bumped).
Since no $2n/3$-quorum on a single target formed, the plurality count is $< 2n/3$, so $|\{i : \sigma_{\mathrm{pre}}.\jtargets[i] \neq T^* \vee \sigma_{\mathrm{pre}}.\jtargets[i] = \bot\}| > n/3$.

\emph{Layer~2 finalize.} Fires when a pending-finality gap is open at the start of the slot ($\sigma_{\mathrm{pre}}.F \prec \sigma_{\mathrm{pre}}.J$) and $F$ did not advance; every $i \notin \sigma_{\mathrm{pre}}.P$ is bumped.
The finality branch of \processHeight, when it fires, sets $\sigma_{\mathrm{post}}.F = \sigma_{\mathrm{pre}}.J \succ \sigma_{\mathrm{pre}}.F$, so this guard is skipped automatically.

\paragraph{Canonical-freshness.}
A vote~$v$ processed at the current slot on canonical chain~$X$ at state-height $h_c = \sigma_{\mathrm{post}}.h$ sets $\sigma.\voteSlots[v.\validator] \neq \bot$ iff $v.\height = h_c$; a \just vote additionally must satisfy $v.\target \preceq X$ and $v.\target.\slot \geq \sigma_{\mathrm{post}}.s_h$ (equivalently $\sigma[v.\target].h = h_c$, \Cref{rem:fresh-equiv}).

\subsection{Leak tightness}

\begin{definition}[Non-finality period]\label{def:non-finality-period}
Fix a chain and write $\sigma_{\mathrm{pre}}^{(t)}, \sigma_{\mathrm{post}}^{(t)}$ for its pre/post-\processHeight snapshots at slot~$t$.
A \emph{non-finality period} is a maximal interval $[t_1, t_T]$ of consecutive slots on which $F$ is stuck, i.e., $\sigma_{\mathrm{post}}^{(t)}.F = \sigma_{\mathrm{pre}}^{(t)}.F$ throughout.
Each such slot falls into exactly one category:
\begin{itemize}
  \item \emph{pending-stuck}: $\sigma_{\mathrm{pre}}^{(t)}.F \prec \sigma_{\mathrm{pre}}^{(t)}.J$ (pending-finality gap open at the start);
  \item \emph{stable-stuck}: $\sigma_{\mathrm{pre}}^{(t)}.F = \sigma_{\mathrm{pre}}^{(t)}.J$ and $J$ does not advance at~$t$;
  \item \emph{transition}: $\sigma_{\mathrm{pre}}^{(t)}.F = \sigma_{\mathrm{pre}}^{(t)}.J$ and the justification branch fires at~$t$ (opening a gap by slot's end).
\end{itemize}
\end{definition}

\begin{theorem}[Non-finality period penalty bound]\label{thm:tightness}
Over any non-finality period of length $T \geq 1$ on any chain, the total penalty increment across all validators assessed by \processInactivityPenalties on that chain is at least $(T - 1) \cdot n/3$.
\end{theorem}

\begin{proof}
Fix a non-finality period $[t_1, t_T]$.
We show (i) at most one slot is transition, and (ii) every pending-stuck or stable-stuck slot contributes $> n/3$ to the total penalty increment.

\emph{(i) At most one transition slot.}
After a transition at~$t$, $\sigma_{\mathrm{post}}^{(t)}.J \succ \sigma_{\mathrm{pre}}^{(t)}.F$; since $F$ does not advance in the period and $J$ is monotone across slots, $\sigma_{\mathrm{pre}}^{(t')}.F \prec \sigma_{\mathrm{pre}}^{(t')}.J$ at every later $t' > t$, making every such $t'$ pending-stuck.

\emph{(ii) Per-slot charge $> n/3$.}

\emph{Pending-stuck.} The Layer~2 finalize guard fires at~$t$.
The finality branch of \processHeight did not fire (else $F$ would advance), so $|\sigma_{\mathrm{pre}}^{(t)}.P| < 2n/3$ and $|\{i : i \notin \sigma_{\mathrm{pre}}^{(t)}.P\}| > n/3$.

\emph{Stable-stuck.} The Layer~2 target guard fires at~$t$.
The justification branch did not fire, so the plurality count on $\sigma_{\mathrm{pre}}^{(t)}.\jtargets$ is $< 2n/3$ (counting $T^* = \bot$ as~$0$); hence the bumped set has size $> n/3$.

At most one slot is transition, so at least $T - 1$ slots are pending-stuck or stable-stuck; each contributes $> n/3$ by~(ii), and Layer~1 adds only non-negative charges.
Summing gives total increment $\geq (T-1) \cdot n/3$.
\end{proof}

\subsection{Honest voting rule}

We state the honest rule explicitly, as it underpins the leak fairness proof.

\begin{definition}[Honest voting rule]\label{def:voting-rule}
Assume the store $\Sigma$ carries a field $\Sigma.V$, the set of votes~$i$ has previously cast.
At each voting opportunity, $i$ constructs its vote by $\constructVote(\Sigma,\, \Omega,\, k)$ where $\Omega$ is~$i$'s current auxiliary state (\Cref{alg:store}) and $k$ is the current slot, defined in \Cref{alg:voting-rule}.
The first-vote check in the (R1) branch ensures at most one \just vote per height~$h_c$, and the target-selection clause preserves E1-safety under view shifts.
\end{definition}

\begin{figure}[H]
\caption{Honest voting rule}\label{alg:voting-rule}
\begin{center}
\begin{boxedminipage}{\textwidth}
\begin{algorithmic}[1]

\Function{\constructVote}{$\Sigma,\; \Omega,\; k$} \Comment{$k$ is the current slot}
  \State $T \gets \getConfirmed(\Sigma,\; \Omega)$
  \State $h_c \gets \sigma[T].h$; \;\; $h_j \gets \sigma[T].h_j$; \;\; $J \gets \sigma[T].J$
  \vspace{3pt}
  \State \textbf{-- Target selection --}
  \If{$\exists\, v \in \Sigma.V$ with $v.\finalizeHeight = h_c$ and $v.\finalizeTarget \neq \bot$}
    \State $T_c \gets v.\finalizeTarget$ \Comment{$i$ is locked at $h_c$}
  \Else
    \State $T_c \gets T$ \label{line:unlocked-target}
  \EndIf
  \vspace{3pt}
  \State \textbf{-- Vote kind and finalize fields --}
  \If{$\not \exists v \in \Sigma.V$ s.t. $v.\height = h_c$} \Comment{first vote at $h_c$: cast \just (R1)}
    \If{every $v \in \Sigma.V$ with $v.\height = h_j$ has $v.\target = J$}
      \State $(f_h,\, f_t) \gets (h_j,\, J)$
    \Else
      \State $(f_h,\, f_t) \gets (\bot,\, \bot)$
    \EndIf
    \State \Return $(i,\; h_c,\; k,\; T_c,\; \just,\; f_h,\; f_t)$
  \Else \Comment{already voted at $h_c$: cast \notar (R2)}
    \State \Return $(i,\; h_c,\; k,\; T_c,\; \notar,\; \bot,\; \bot)$
  \EndIf
\EndFunction

\end{algorithmic}
\end{boxedminipage}
\end{center}
\end{figure}

Only the (R1) branch produces a \just vote and, gated by the first-vote check, fires at most once per~$h_c$, ensuring per-height \just-uniqueness.
The (R2) branch is the \emph{resubmission} fallback: fired at every subsequent voting opportunity while canonical stays at~$h_c$.

Checking \emph{every} prior vote of~$i$ at~$h_j$ (not just the (R1) one) is required because a (R2) at~$h_j$ may have chosen a target differing from the (R1) target if~$i$'s view shifted; committing $f_t = J$ at height~$h_j$ requires (by E1) that every $h_j$-vote of~$i$ have target~$J$, so in case of a split the conditional fires $(\bot, \bot)$: $i$ misses the finalization at~$h_j$ but remains E1-safe.

\begin{lemma}[Honest E1 safety]\label{lem:honest-e1-safety}
If validator~$i$ follows \Cref{def:voting-rule}, then $i$ never produces a vote-pair $(v_1, v_2)$ with $v_2.\finalizeTarget \neq \bot$, $v_1.\height = v_2.\finalizeHeight$, and $v_1.\target \neq v_2.\finalizeTarget$.
\end{lemma}

\begin{proof}
Write $h_\ast \coloneqq v_2.\finalizeHeight$ and $T_\ast \coloneqq v_2.\finalizeTarget \neq \bot$.
Only (R1) produces finalize-bearing votes, so $v_2$ is an (R1) vote with $(f_h, f_t) = (h_j, J) = (h_\ast, T_\ast)$ read from the canonical state at~$v_2$'s firing.

\emph{Case~A ($v_1$ cast before~$v_2$).}
The (R1) finalize-field conditional at~$v_2$ attached $(h_\ast, T_\ast)$ only because every prior vote of~$i$ at height~$h_\ast$ had target~$T_\ast$; in particular $v_1.\target = T_\ast$.

\emph{Case~B ($v_1$ cast after~$v_2$).}
At $v_1$'s firing, the voting height is $h_c = v_1.\height = h_\ast$, and~$i$ has the prior vote~$v_2$ with $\finalizeHeight = h_\ast = h_c$ and $\finalizeTarget = T_\ast \neq \bot$: $i$ is \emph{locked at $h_c$}, and the target-selection clause of \Cref{def:voting-rule} forces $T_c = T_\ast$ regardless of~$i$'s canonical view at that moment.
Hence $v_1.\target = T_c = T_\ast$.

\emph{Case~C ($v_1 = v_2$).}
Vacuous: $v_2.\height = h_c \neq h_\ast = v_2.\finalizeHeight$ (every height advance sets $h_j \gets h$ before $h \gets h + 1$, so $h_j < h$ in any post-advance state).
\end{proof}

\subsection{Leak fairness}

Fairness is local to a single honest validator~$i$: we fix~$i$, its local store~$\Sigma$, its auxiliary state~$\Omega$, and the canonical chain $X = \getConfirmed(\Sigma, \Omega)$.
The claim is that the vote~$v$ produced by \Cref{def:voting-rule}, once included on any extension of~$X$ still at the same state-height, exempts~$i$ from Layer~1 penalties at that state-height.
Height advances reset $\jtargets$ and $\voteSlots$ (\Cref{alg:state-machine}), so Layer-2 state at the next height is independent of~$i$'s prior votes and no ``lock'' issue propagates.

\begin{lemma}[Per-height justification uniqueness]\label{lem:just-unique-height}
If $f < n/3$ and all honest validators follow \Cref{def:voting-rule}: for any two blocks $C_1, C_2$ justified at the same height~$h$ on any chains processed by the node, $C_1 = C_2$.
\end{lemma}

\begin{proof}
Each justification at height~$h$ on a chain yields a \just-quorum $Q_k$ ($k \in \{1, 2\}$) of size $\geq 2n/3$ signing a \just vote at height~$h$ for $C_k$; quorum intersection gives $|Q_1 \cap Q_2| \geq n/3$.
Since $f < n/3$, some $v \in Q_1 \cap Q_2$ is honest; by \Cref{def:voting-rule}, only (R1) emits a \just vote and it fires at most once per height, so~$v$'s single \just vote at~$h$ has both $C_1$ and $C_2$ as target, forcing $C_1 = C_2$.
\end{proof}

\begin{lemma}[Finalize vote implies justification]\label{lem:no-finalize-without-just}
If~$i$ has cast a vote~$v$ with $v.\finalizeHeight = h'$ and $v.\finalizeTarget = T_f \neq \bot$, then~$T_f$ was justified at height~$h'$ on some processed chain by the time~$v$ was cast.
\end{lemma}

\begin{proof}
Only (R1) produces $\finalizeHeight \neq \bot$, so~$v$ came from (R1) at an earlier moment with~$i$'s local store~$\Sigma_0$ and auxiliary state~$\Omega_0$; by the (R1) finalize-field conditional, $(f_h, f_t) = (h_j, J) = (h', T_f)$ read from~$i$'s then-canonical $X_0 = \getConfirmed(\Sigma_0, \Omega_0)$, so $\sigma[X_0].h_j = h'$ and $\sigma[X_0].J = T_f$.
A state has $(h_j, J) = (h', T_f)$ only after the justification branch of \processHeight fired on that chain with target~$T_f$ at pre-state height~$h'$, i.e.\ some ancestor $B' \preceq X_0$ triggered justification of~$T_f$ at height~$h'$.
\end{proof}

\begin{lemma}[Canonical-fresh honest vote]\label{lem:canonical-fresh-honest-vote}
Unless $\geq n/3$ validators are slashable, and assuming $f < n/3$ and all honest validators follow \Cref{def:voting-rule}: let~$i$ be honest with local store~$\Sigma$ and auxiliary state~$\Omega$, $X = \getConfirmed(\Sigma, \Omega)$, and $h_c = \sigma[X].h$.
Let~$v$ be the vote~$i$ produces at the current slot.
Then $v.\target \preceq X$ with $\sigma[v.\target].h = h_c$.
In particular, $v$ is canonical-fresh on every extension $X' \succeq X$ with $\sigma[X'].h = h_c$.
\end{lemma}

\begin{proof}
By \Cref{def:voting-rule}, $v.\height = h_c$ and $v.\target = T_c$, where $T_c$ is either an ancestor of~$X$ at state-height~$h_c$ (unlocked case, satisfying the claim directly), or a lock target~$L$ from a prior (R1) vote of~$i$ with $\finalizeHeight = h_c$ (locked case).

\emph{Locked case.}
The prior (R1) was cast at an earlier moment when~$i$'s then-canonical had $h_j = h_c$, so a justification at height~$h_c$ fired on some processed chain at that moment.
\Cref{lem:no-finalize-without-just} gives the descriptor $(h_c, L)$ offered to \updateForkChoiceRoot at that time, and the filter $\Sigma.F \preceq L$ passes under $f < n/3$ (\Cref{rem:fs-invariant} makes $\Sigma.F$ finalized or genesis; \Cref{lem:certchain} with $h_c \geq h_{F}$ gives $\Sigma.F \preceq L$).
Hence \updateForkChoiceRoot raised $\sigma[\Sigma.J].h$ to $\geq h_c$, and \Cref{lem:J-monotone} preserves this bound: currently $\sigma[\Sigma.J].h \geq h_c$.

$X$ descends from~$R$ (\Cref{alg:store}), so state-height monotonicity (\Cref{lem:slotmono}) gives $h_c = \sigma[X].h \geq \sigma[R].h$.
If $R = \Sigma.J$, then $\sigma[\Sigma.J].h \leq h_c$ and the previous bound force $\sigma[\Sigma.J].h = h_c$; \Cref{lem:just-unique-height} then forces $\Sigma.J = L$, so $R = L$.
If $R$ is the fallback (a block at state-height $h^\star - 1$), $\sigma[R].h = h^\star - 1 \leq h_c$ together with the fallback-branch trigger $\sigma[\Sigma.J].h < h^\star - 1$ would give $\sigma[\Sigma.J].h < h_c$, contradicting the bound above.
Hence $R = L$ and $T_c = L \preceq X$ with $\sigma[L].h = h_c$.
\end{proof}

\begin{theorem}[Layer~1 leak fairness]\label{thm:fairness}
Unless $\geq n/3$ validators are slashable, and assuming $f < n/3$ and all honest validators follow \Cref{def:voting-rule}: let~$i$ be an honest validator, $\Sigma$ its local store, $\Omega$ its auxiliary state, $X = \getConfirmed(\Sigma, \Omega)$ its canonical chain, and $h_c = \Sigma.\sigma[X].h$ the current canonical state-height.
Let $v$ be the vote~$i$ produces at the current slot.
Let $\Sigma'$ be any store obtained by processing additional blocks and votes on top of~$\Sigma$, with auxiliary state~$\Omega'$, such that $X' = \getConfirmed(\Sigma', \Omega')$ is an extension of~$X$ with $\Sigma'.\sigma[X'].h = h_c$, and suppose~$v$ is included on~$X'$.
Then for every slot processed on~$X'$ from that point onward while canonical state-height stays at~$h_c$, the Layer~1 increment to $\sigma.\penalties[i]$ on canonical is zero.
\end{theorem}

\begin{proof}
\Cref{lem:canonical-fresh-honest-vote} gives $T_c \coloneqq v.\target \preceq X \preceq X'$ with $\sigma[T_c].h = h_c$.
The hypothesis $\Sigma'.\sigma[X'].h = h_c$ with \Cref{lem:slotmono} gives $\Sigma'.\sigma[X'].s_h = \Sigma.\sigma[X].s_h$; \Cref{rem:fresh-equiv} then gives $T_c.\slot \geq \Sigma'.\sigma[X'].s_h$.
So~$v$ is canonical-fresh on~$X'$ (passing target freshness in the \just case; height freshness alone in the \notar case), and processing~$v$ sets $\voteSlots[i] \neq \bot$.
$\voteSlots$ resets only at a height advance, which does not fire while state-height stays at~$h_c$, so $\voteSlots[i] \neq \bot$ persists and Layer~1 does not bump~$i$.
\end{proof}

The $f < n/3$ hypothesis is used only in \Cref{lem:canonical-fresh-honest-vote} via \Cref{lem:just-unique-height}, to collapse the lock target~$L$ onto $R$ in the $h_c = h_R$ case.

\end{document}
